
% Hoe wil ik dat de thesis eruit ziet:

% 1. Introductie - belang corners en relevante
% 2. Problem statement, erg uitgebreid 
% 3. Academic Context (Literature Review)
% 4. Methodology \& Data zoals nu
% 5. Results
% 6. Use cases / Application of results
% 7. Discussion
% 8. Conclusions



\chapter{Methods}
\label{sec:methods}

\section{Research Design}
\label{sec:research_design}

This study adopts a quantitative, data-driven research design that sequentially
builds five analytical layers, each corresponding to one sub-question (SQ).
The overarching logic is as follows: corner kicks are first characterised from both
sides of the ball (SQ1: defensive setup; SQ2: offensive pattern), then matched to
quantify effectiveness (SQ3), enriched with individual player attributes (SQ4), and
finally synthesised into a practical preparation tool (SQ5).
All analyses are applied to the same dataset, ensuring consistency across sub-questions.
Figure~\ref{fig:pipeline} provides an overview of the analytical pipeline.

% \begin{figure}[h]
%   \centering
%   \includegraphics[width=\linewidth]{figures/pipeline_overview.pdf}
%   \caption{Overview of the five-layer analytical pipeline.
%            Arrows indicate data flow between layers.}
%   \label{fig:pipeline}
% \end{figure}

\section{Data}
\label{sec:data}

\subsection{Data Sources}
\label{sec:data_sources}

All data were provided by SciSports through FC Den Bosch. Two data types were used.

\subsubsection{Event data.}
Event data contain one record per on-ball action, including the action type
(e.g.\ \texttt{CORNER\_CROSSED}, \texttt{CORNER\_SHORT}, \texttt{HEADER},
\texttt{CLEARANCE}), the player and team involved, the ball's start and end
coordinates in metres (origin at pitch centre, $x$-axis along the length,
$y$-axis along the width), the timestamp in milliseconds, and derived metrics
such as delivery length, angle, expected goals (xG), and expected assists (xA).
Additionally, each event contains a \texttt{sequenceId} that groups all actions
within a single possession sequence, enabling reconstruction of what happened
after every corner kick.

\subsubsection{Tracking data.}
Tracking data record the $(x, y)$ position of every player and the ball at
10~Hz (one frame per 100~ms). Player records include a unique player identifier
and shirt number. Ball records include an additional $z$-coordinate (height in
metres) allowing computation of ball flight characteristics.
Tracking data are provided in positional JSON files aligned to the same
millisecond clock as the event data. Where a corner event is present in the
event data and a corresponding tracking frame exists within a 200~ms tolerance,
the corner is marked as \textit{synchronised} (\texttt{synced = True}).

\section{Dataset Composition}
\label{sec:dataset_composition}

Data span the 2025–2026 Dutch Eerste Divisie season, covering 281 unique matches
from 20 clubs. All 1,893 corner kicks that appear in the event data were extracted,
of which all 1,893 are synchronised with tracking data. These consist of 1,401
direct (crossed) corners and 492 short corners.
Table~\ref{tab:dataset_overview} summarises the dataset.

\begin{table}[h]
  \centering
  \caption{Dataset summary.}
  \label{tab:dataset_overview}
  \begin{tabular}{lr}
    \toprule
    \textbf{Property} & \textbf{Value} \\
    \midrule
    Matches & 281 \\
    Clubs & 20 \\
    Total corner kicks & 1,893 \\
    \quad Direct (crossed) & 1,401 \\
    \quad Short & 492 \\
    Synced with tracking & 1,893 (100\%) \\
    Manually labelled corners & 86 \\
    \quad Labelled attacker rows & 389 \\
    \quad Labelled defender rows & 509 \\
    Tracking frame rate & 10 Hz \\
    \bottomrule
  \end{tabular}
\end{table}

\subsection{Manual Labelling}
\label{sec:manual_labelling}

A subset of 86 corners was manually labelled by the primary researcher using a
custom labelling tool (\texttt{labeling\_tool.py}) that overlays tracking data
on a 2D pitch visualisation. For each corner, the following were annotated:
\begin{itemize}
  \item \textbf{Delivery type}: swing direction and target zone (see
        Section~\ref{sec:sq2_delivery});
  \item \textbf{Ball flight}: one of five trajectory categories
        (see Section~\ref{sec:sq2_flight});
  \item \textbf{Attacker run roles}: per-player tactical role for all attacking
        outfield players (see Section~\ref{sec:sq2_runs});
  \item \textbf{Defender roles}: per-player defensive role
        (see Section~\ref{sec:sq1});
  \item \textbf{Man-marking assignments}: for defenders classified as man-markers,
        the jersey number of the specific attacker being tracked.
\end{itemize}
Labels were recorded in a structured spreadsheet
(\texttt{labeling\_sheet\_v2.csv}) and used as ground truth for model training
and validation.
Inter-rater reliability was not formally assessed, as only one annotator was
used; this is acknowledged as a limitation in Section~\ref{sec:limitations}.

\paragraph{Zone system.}
The pitch area surrounding the goal was divided according to the zone system
depicted in Figure~\ref{fig:zone_system}. A central column of width 18.32~m
(the goal area width) was extended perpendicularly from the goal line up to
the edge of the penalty arc ($\approx20.15$~m from the goal line). This column
was subdivided into three lateral lanes of equal width (6.11~m each): column~1
(near post), column~2 (central, between the posts), and column~3 (far post).
Four depth bands were defined within this column:
\begin{itemize}
  \item \textbf{GA} (Goal Area): 0–5.5~m from the goal line ($x > 47.0$~m);
  \item \textbf{PA} (Penalty Area): 5.5–11.0~m ($x > 41.5$~m);
  \item \textbf{PE} (Penalty Edge): 11.0–16.5~m ($x > 36.0$~m);
  \item \textbf{PC} (Penalty Curve/Arc): 16.5–20.15~m ($x > 32.35$~m).
\end{itemize}
Beyond the central column, two lateral regions were defined: \textbf{Front Zone}
(near side of the penalty area, $y_\mathrm{near} > 9.16$~m and $x > 36.0$~m)
and \textbf{Edge Zone} (far side, $y_\mathrm{near} < -9.16$~m and
$x > 36.0$~m), where $y_\mathrm{near}$ denotes the signed lateral distance from
the goal centre after normalisation to attacking-right coordinates.

% \begin{figure}[h]
%   \centering
%   \includegraphics[width=0.7\linewidth]{figures/zone_system.pdf}
%   \caption{Zone system used for delivery and player position classification.
%            The central column (18.32~m wide) is divided into four depth bands
%            and three lateral lanes. GA = Goal Area, PA = Penalty Area,
%            PE = Penalty Edge, PC = Penalty Curve. Numbers 1--3 denote near,
%            central, and far lateral columns respectively.}
%   \label{fig:zone_system}
% \end{figure}

\section{SQ1: Defensive Corner Setup Classification}
\label{sec:sq1}

\subsection{Individual Defender Role Classification}
\label{sec:sq1_individual}

The classification of individual defender roles was developed in a companion study
(\emph{Defensive Corners Onderzoek}) and is described in full detail therein.
Briefly, an ensemble classifier (Random Forest, Extra Trees, and Gradient Boosting,
averaged) assigns each defending outfield player to one of four roles:
\texttt{MAN} (man-marker), \texttt{ZONAL} (covers a pitch zone),
\texttt{SHORT} (guards the short-corner option), or \texttt{COUNTER}
(positioned for transition defence).
The classifier uses pair-wise features computed over a 5,000~ms and a 3,000~ms
tracking window prior to the corner kick, including mean pairwise distance
to each attacker, motion correlation, tracking fraction, distance to the ball,
and mutual nearest-neighbour indicators.
A leave-one-match-out cross-validation on the manually labelled dataset achieved
a role accuracy of 93\% and a marking assignment accuracy (Hungarian-algorithm
assignment of man-markers to specific attackers) of approximately 80\%.

\subsection{Team-Level Defensive Strategy}
\label{sec:sq1_team}

Individual role labels are aggregated per corner into a team-level defensive
strategy. Three strategy dimensions are defined:
\begin{enumerate}
  \item \textbf{Marking intensity}: the number of man-markers present
        (0–1: predominantly zonal; 2–3: mixed; $\geq$4: predominantly man-marking);
  \item \textbf{Short defender present}: binary (yes/no), indicating whether at
        least one defender is assigned to guard the short-corner option;
  \item \textbf{Counter defender present}: binary (yes/no), indicating whether at
        least one defender remains in the defensive half.
\end{enumerate}
These dimensions produce a compact categorical label for each corner, referred
to as the \textit{defensive strategy type}. This label serves as the defensive
predictor variable in SQ3.

\subsection{Team-Level Profile Estimation}
\label{sec:sq1_team_profile}

For each club with a minimum of 10 home corners in the dataset, a \textit{defensive
profile} is computed as the distribution over defensive strategy types across all
their corners as the defending team. This profile is used in SQ5 (the tactical
preparation tool) to characterise how a specific upcoming opponent typically
defends corners.

\section{SQ2: Offensive Corner Pattern Classification}
\label{sec:sq2}

Offensive corner patterns are characterised along three independent dimensions:
(i) delivery type, (ii) ball flight trajectory, and (iii) attacker movement
patterns. Together, these define the \textit{offensive corner profile} for each
corner.

\subsection{Delivery Type}
\label{sec:sq2_delivery}

Delivery type is defined as the combination of swing direction and target zone.

\subsubsection{Swing direction.}
Swing direction is determined from the kicker's foot and the corner side using
the rule: a right-footed kick from the right corner flag produces an inswing
(\texttt{IN}); the same foot from the left corner flag produces an outswing
(\texttt{OUT}). Left foot follows the mirror rule. Where foot data are ambiguous
(\texttt{FEET}), the swing label provided by the SciSports event pipeline is
used. Swing direction is validated against the manually labelled subset, yielding
92\% agreement.

\subsubsection{Target zone.}
For direct corners, the target zone is derived from the ball's endpoint
coordinates $(x_\mathrm{end}, y_\mathrm{end})$ in the event data using the zone
classification described in Section~\ref{sec:data}. Validation against the 51~manually labelled delivery zones yields the following
agreement rates: full zone label (e.g.\ \texttt{GA2}): 64\%;
depth band only (GA/PA/PE/PC): 74\%; lateral column only (near/central/far):
76\%.
The remaining discrepancies are attributable to the fact that the SciSports
event data records the ball's position at the moment of \emph{first player
contact} (header, clearance) rather than the ball's intended landing zone.
An extensive comparison of alternative tracking-based methods
(ball position at minimum forward distance, at mid-flight time, and at peak
height) found none to exceed 26\% zone accuracy, confirming that the event
endpoint is the best available automated proxy.
Consequently, the SQ3 effectiveness analysis uses the full zone label only for
the 51~manually labelled corners, and relies on the more reliable depth band
(74\% accuracy) and swing direction (92\% accuracy) dimensions for the full
1,893-corner dataset. The full zone label for unlabelled corners should be
treated as approximate and is not used in inferential models.

The full set of delivery type codes is given in Table~\ref{tab:delivery_types}.

\begin{table}[h]
  \centering
  \caption{Delivery type codes for direct corners. \texttt{in}/\texttt{out}
           denote inswing and outswing respectively. Rows marked with $*$ are
           outer-zone deliveries.}
  \label{tab:delivery_types}
  \begin{tabular}{ll|ll}
    \toprule
    \textbf{Code} & \textbf{Description} &
    \textbf{Code} & \textbf{Description} \\
    \midrule
    \texttt{inGA1} & Inswing, near GA  & \texttt{outGA1} & Outswing, near GA \\
    \texttt{inGA2} & Inswing, central GA & \texttt{outGA2} & Outswing, central GA \\
    \texttt{inGA3} & Inswing, far GA   & \texttt{outGA3} & Outswing, far GA \\
    \texttt{inPA1} & Inswing, near PA  & \texttt{outPA1} & Outswing, near PA \\
    \texttt{inPA2} & Inswing, central PA & \texttt{outPA2} & Outswing, central PA \\
    \texttt{inPA3} & Inswing, far PA   & \texttt{outPA3} & Outswing, far PA \\
    \texttt{inPS}  & Inswing, penalty spot area & \texttt{outPS} & Outswing, penalty spot \\
    \texttt{inPE}  & Inswing, penalty edge & \texttt{outPE} & Outswing, penalty edge \\
    \texttt{inFront}$^*$  & Inswing, near-side outer &
    \texttt{outFront}$^*$ & Outswing, near-side outer \\
    \texttt{inEdge}$^*$   & Inswing, far-side outer &
    \texttt{outEdge}$^*$  & Outswing, far-side outer \\
    \texttt{inOther}$^*$  & Inswing, outside PA depth &
    \texttt{outOther}$^*$ & Outswing, outside PA depth \\
    \midrule
    \multicolumn{4}{l}{\texttt{SHORT\_1PASS}, \texttt{SHORT\_2PASS},
                        \texttt{SHORT\_SHOT}, \texttt{SHORT\_RECYCLE}
                        (short corners, see Section~\ref{sec:sq2_short})} \\
    \bottomrule
  \end{tabular}
\end{table}

\subsection{Ball Flight Trajectory}
\label{sec:sq2_flight}

Ball flight is classified from the peak height ($z_\mathrm{peak}$, metres) of
the ball during the delivery phase, extracted from tracking data. Four categories
are defined. An initial five-category scheme including a separate
\texttt{DRIVEN} band was tested but subsequently merged into \texttt{MEDIUM}
after manual inspection confirmed that the 3.0–5.5~m range did not represent a
meaningfully distinct tactical category from the broader 3.0–8.0~m arc:

\begin{table}[h]
  \centering
  \caption{Ball flight categories and their $z_\mathrm{peak}$ thresholds.}
  \label{tab:flight}
  \begin{tabular}{lll}
    \toprule
    \textbf{Category} & \textbf{Threshold} & \textbf{Description} \\
    \midrule
    \texttt{GROUND}  & $z_\mathrm{peak} \leq 0.5$~m & Along the ground; typical of cutbacks \\
    \texttt{LOW}     & $0.5 < z_\mathrm{peak} \leq 3.0$~m & Low-driven, flat delivery \\
    \texttt{MEDIUM}  & $3.0 < z_\mathrm{peak} \leq 8.0$~m & Flat-to-arc delivery \\
    \texttt{FLOATED} & $z_\mathrm{peak} > 8.0$~m & High arc, slow descent \\
    \bottomrule
  \end{tabular}
\end{table}

\subsection{Short Corner Sub-types}
\label{sec:sq2_short}

Short corners are classified using a detailed taxonomy derived from 21~manually
labelled short corners, which records three components: (i) the number of passes
in the sequence (counting the corner delivery itself as pass~1), (ii) the number
of distinct attacking players involved, and (iii) the outcome of the sequence.
Labels follow the format:
\[
  \texttt{SHORT\_\{n\}pass\{k\}players-\{result\}}
\]
where $n$ is the total pass count, $k$ the number of attacking players involved,
and \textit{result} $\in$ \{\texttt{cross}, \texttt{shot}, \texttt{ballloss}\}.
For example, \texttt{SHORT\_2pass3players-cross} describes a short corner in
which the delivery is followed by one additional pass involving three attackers
before a cross is delivered.

This classification is applied to all 492 short corners in the dataset by
tracing the event sequence from the corner delivery until either a cross, shot,
or defensive ball-recovery event occurs. Only genuine PASS events (not carries
or dribbles) are counted towards the pass total. An automated classifier
achieves 81\% exact-label accuracy against the 21 ground-truth labels, with
result accuracy (cross/shot/ballloss) of 86\% and player-count accuracy of 86\%.
The four remaining discrepancies involve edge cases (crosses blocked into an
out-of-play event, or sequences spanning multiple event sequence identifiers).

For the SQ3 effectiveness analysis, short corner labels are further aggregated
into nine bucketed types based on pass count ($\leq$1, 2, $\geq$3) and outcome
(cross, shot, ballloss), ensuring a minimum cell size of $n \geq 20$
per analytical group.

\subsection{Attacker Movement Patterns}
\label{sec:sq2_runs}

Player movement during the 3-second pre-delivery window is characterised at two
levels: the \textit{individual running line} and the \textit{corner movement
profile}.

\paragraph{Individual running line.}
Each attacking outfield player (excluding the corner taker) is assigned a
geometric running line label based on their starting zone (3~s before delivery)
and ending zone (2~s after delivery), as well as their displacement:

\begin{itemize}
  \item \texttt{STATIC}: displacement $< 2.5$~m; player does not make a meaningful run;
  \item \texttt{NEAR\_POST}: player ends in GA1 or PA1 (near-post side);
  \item \texttt{FAR\_POST}: player ends in GA3 or PA3 (far-post side);
  \item \texttt{CENTRAL}: player ends in GA2 or PA2 (central lane);
  \item \texttt{CROSSING\_FN}: player starts in a far-side zone and ends in a
        near-side zone (far-to-near crossing run);
  \item \texttt{CROSSING\_NF}: player starts in a near-side zone and ends in a
        far-side zone (near-to-far crossing run);
  \item \texttt{EDGE}: player ends in the PE or PC depth band (positions for
        second-phase shooting);
  \item \texttt{SECOND\_PHASE}: player ends in an ED or OUT zone (outside the
        penalty area; second-ball positioning).
\end{itemize}

Running line labels are derived rule-based from zone transitions and yield 99.5\%
leave-one-corner-out accuracy on the manually labelled subset.

\subsubsection{Corner movement profile.}
The distribution of individual running lines within a corner is summarised as a
\textit{movement profile string}: an unordered count-based label recording how
many players execute each running line type. For example,
\texttt{2NEAR\_1CENTRAL\_3SECOND\_PHASE\_4STATIC} denotes a corner in which
two players run to the near post, one to the central zone, three position for
second-phase balls, and four remain static. This label serves two purposes:
first, it characterises the offensive setup in full detail; second, two
categorical flags are derived from it:
\begin{itemize}
  \item \textbf{Static vs.\ dynamic}: a corner is classified as \textit{static}
        if fewer than three players make a displacement $\geq 2.5$~m, and as
        \textit{dynamic} otherwise;
  \item \textbf{Dominant run direction}: the running line type with the highest
        count (excluding \texttt{STATIC} and \texttt{SECOND\_PHASE}), used as a
        single summary label when the full profile would produce overly sparse
        analytical cells.
\end{itemize}

\subsubsection{Attacker tactical role (supplementary).}
For the 86 manually labelled corners only, each attacker is also assigned a
\textit{tactical role}: \texttt{TARGET} (intends to receive the ball),
\texttt{DECOY} (creates space by pulling defenders), \texttt{BLOCK\_GK}
(screens the goalkeeper), \texttt{BLOCK\_DEF} (screens a specific defender),
\texttt{SECOND\_BALL} (positions for rebounds), or \texttt{STATIC}. A Random
Forest classifier trained on this subset achieves 63\% accuracy
(leave-one-corner-out). Because TARGET and DECOY are subjectively defined by
tactical intent rather than observable geometry, this model is reported as a
supplementary finding; the geometric running line classification is used as the
primary offensive movement descriptor in SQ3.

\subsection{Offensive Corner Profile}
\label{sec:sq2_profile}

The three dimensions — delivery type, ball flight category, and corner movement
profile — together constitute the \textit{offensive corner profile}. Due to the
high granularity of the full cross-product (potentially approaching one unique
label per corner), the SQ3 effectiveness analysis operates on individual dimensions
and their pairwise combinations rather than the full cross-product profile.
Table~\ref{tab:offensive_dims} lists the analytical dimensions and their levels.

\begin{table}[h]
  \centering
  \caption{Analytical dimensions of the offensive corner profile.
           Dimensions are used individually or in pairwise combinations
           in SQ3.}
  \label{tab:offensive_dims}
  \begin{tabular}{llr}
    \toprule
    \textbf{Dimension} & \textbf{Levels} & \textbf{$N$ levels} \\
    \midrule
    Swing direction & \texttt{IN}, \texttt{OUT} & 2 \\
    Target zone depth & \texttt{GA}, \texttt{PA}, \texttt{PE/PC}, \texttt{FRONT/EDGE} & 4 \\
    Lateral target column & \texttt{NEAR} (1), \texttt{CENTRAL} (2), \texttt{FAR} (3) & 3 \\
    Ball flight & \texttt{GROUND}, \texttt{LOW}, \texttt{MEDIUM}, \texttt{FLOATED} & 4 \\
    Setup type & \texttt{STATIC}, \texttt{DYNAMIC} & 2 \\
    Dominant run direction & \texttt{NEAR\_POST}, \texttt{FAR\_POST},
                             \texttt{CENTRAL}, \texttt{CROSSING}, \texttt{EDGE} & 5 \\
    Short vs.\ direct & \texttt{SHORT}, \texttt{DIRECT} & 2 \\
    \bottomrule
  \end{tabular}
\end{table}

\section{SQ3: Effectiveness of Offensive Strategies Against Defensive Structures}
\label{sec:sq3}

\subsection{Outcome Variable}
\label{sec:sq3_outcome}

Following the recommendation in recent systematic reviews
\citep[e.g.][]{Plakias2025, Casal2015}, corner kick effectiveness is
measured using two complementary outcome indicators.

\subsubsection{Corner Threat Score (CTS).}
\citet{Lamberts2023} proposes the Corner Threat Score, defined as the harmonic
mean of the delivery quality (xA) and conversion quality (xG) components:
\begin{equation}
  \mathrm{CTS} = \frac{2 \cdot \overline{\mathrm{xG}} \cdot \overline{\mathrm{xA}}}
                       {\overline{\mathrm{xG}} + \overline{\mathrm{xA}}},
  \label{eq:cts}
\end{equation}
where $\overline{\mathrm{xA}}$ is the SciSports expected-assist value on the
corner delivery event, and $\overline{\mathrm{xG}}$ is the maximum expected-goals
value of any shot attempt in the same possession sequence after the corner.
When no shot occurs in the sequence, $\overline{\mathrm{xG}} = 0$ and CTS is
set to zero, which is conservative but avoids imputation.
CTS is the \textit{primary continuous effectiveness measure}.

In this study, $\overline{\mathrm{xG}}$ is extracted by tracing each corner's
\texttt{sequenceId} through the raw match event files and selecting the highest
xG value among all subsequent \texttt{SHOT}, \texttt{HEADER}, and
\texttt{SHOT\_FREE\_KICK} events in that sequence.
Across all 1,401 direct corners, 347 (24.8\%) generated at least one shot;
the remaining corners contribute CTS~= 0.
The overall mean CTS for direct corners is 0.0175
($\overline{\mathrm{xA}} = 0.018$; $\overline{\mathrm{xG}} = 0.040$).

\subsubsection{Shot generation.}
A binary indicator ($y_1 \in \{0,1\}$) of whether any shot event occurs within
the corner sequence. Shot generation reflects whether the corner created
a scoring opportunity independent of shot quality, and is used alongside CTS
in the regression analysis.

\subsection{Analytical Strategy}
\label{sec:sq3_analysis}

The effectiveness analysis proceeds in three steps.

\subsubsection{Step 1 — Descriptive cross-tabulation.}
For each combination of offensive dimension and defensive strategy type,
mean CTS, mean xA, mean sequence xG, and shot rate are computed and displayed
as a two-dimensional effectiveness matrix. Fisher's exact test is applied to shot
generation counts for each cell pair; the Bonferroni correction adjusts for
multiple comparisons. Only cells with $n \geq 15$ corners are reported; those
with $n < 15$ are noted as insufficiently powered.

\subsubsection{Step 2 — Logistic regression (shot generation).}
A logistic regression model is fitted with shot generation ($y_1$) as the
outcome:
\begin{equation}
  \log\frac{P(y_1=1)}{P(y_1=0)} = \beta_0 +
    \beta_1 \cdot \mathrm{swing} +
    \beta_2 \cdot \mathrm{zone\_depth} +
    \beta_3 \cdot \mathrm{flight} +
    \beta_4 \cdot \mathrm{setup\_type} +
    \beta_5 \cdot \mathrm{def\_strategy} +
    \beta_6 \cdot (\mathrm{swing} \times \mathrm{def\_strategy}) + \varepsilon,
  \label{eq:logit}
\end{equation}
where categorical predictors are dummy-coded. The interaction term tests whether
the effectiveness of a particular swing direction differs by defensive strategy
type — the core hypothesis of SQ3. Model fit is assessed using the area under
the receiver operating characteristic curve (AUC-ROC) and the log-likelihood
ratio test against an intercept-only baseline.

\subsubsection{Step 3 — Linear regression (CTS).}
An ordinary least squares (OLS) regression of the same predictors on CTS is
estimated alongside Equation~\ref{eq:logit}. Given the zero-inflated
distribution of CTS (75.2\% of direct corners have CTS~= 0 owing to no
subsequent shot), a two-part hurdle model \citep[see][for the approach in
sports contexts]{Rathke2017} is considered as a robustness check.

\subsubsection{Sample Size and Power}
\label{sec:sq3_power}

With 1,401 direct corners, a swing dimension (2 levels) and three defensive
strategy types (3 levels), the average cell size is approximately
$1401 / (2 \times 3) \approx 233$ corners. Interaction cells using the full
delivery type code (22 levels) are considerably smaller; these are therefore
not included in the regression model, and the analysis restricts the delivery
dimension to the coarser decomposition in Table~\ref{tab:offensive_dims}.

\subsection{Descriptive Effectiveness Results}
\label{sec:sq3_descriptive}

\subsubsection{Delivery type.}
Table~\ref{tab:cts_by_delivery} reports mean CTS by delivery type for all
groups with $n \geq 10$ crossed corners.
Central deliveries into the penalty area (IN\_PA3: CTS~= 0.033; OUT\_GA2:
CTS~= 0.032; IN\_GA2: CTS~= 0.029) are substantially more threatening than
near-post and far-post deliveries to the same depth band
(IN\_GA1: CTS~= 0.015; OUT\_GA1: CTS~= 0.004).
Deliveries to the \texttt{FRONT} zone (in front of the six-yard box) and to
the \texttt{EDGE} zone are consistently below average (CTS~$< 0.005$).
In-swing corners (CTS~= 0.018) marginally outperform out-swing corners
(CTS~= 0.017).

\subsubsection{Ball flight.}
Medium-height deliveries ($3$--$8$~m peak height; CTS~= 0.021) generate
more threat than floated deliveries ($>8$~m; CTS~= 0.009) or ground-level
deliveries ($<0.5$~m; CTS~= 0.001). Low deliveries ($0.5$--$3$~m;
CTS~= 0.013) occupy an intermediate position. These patterns suggest that
flatter, more direct deliveries are harder to defend than high, hanging
crosses in the Eerste Divisie context.

\subsubsection{Attacker movement (running lines).}
Running line data are available for 334 of the 1,401 crossed corners
(corners with full tracking coverage). Among corners with a dominant run
direction, central-dominated setups (\texttt{n\_central} highest;
CTS~= 0.034, $n = 13$) and second-phase setups (CTS~= 0.024, $n = 158$)
outperform corners dominated by static positioning (CTS~= 0.013, $n = 125$)
and near-post runs (CTS~= 0.003, $n = 27$). The near-post finding is
noteworthy: near-post runners draw defenders to the near post and may
\textit{reduce} the attacking threat, consistent with the role of these
runners as blockers rather than primary receivers. The majority of corners
($n = 1{,}067$) lack tracking coverage (\texttt{UNKNOWN}); this limits the
power of running line comparisons.

\subsubsection{Attacker tactical roles (86 labelled corners).}
For the manually labelled subset, the presence of blockers is associated
with divergent outcomes: corners featuring a \texttt{BLOCK\_GK} player
(screens the goalkeeper; $n = 14$) show higher mean CTS (0.019) than those
without (0.014), consistent with the blocker creating a direct shooting
lane. By contrast, corners with a \texttt{BLOCK\_DEF} player (blocks a
specific outfield defender; $n = 11$) show lower mean CTS (0.003 vs.\
0.016), possibly reflecting that these corners are executed under conditions
where the defence is concentrated. Corners featuring three decoy runners
(\texttt{DECOY}~= 3; $n = 15$) achieve the highest mean CTS (0.036),
suggesting that systematic decoy movement creates exploitable defensive gaps.

\subsubsection{Defensive strategy.}
Among the 86 labelled corners, corners defended with a predominantly zonal
structure ($n = 37$; CTS~= 0.020) appear more threatening than those with a
predominantly man-marking structure ($n = 6$; CTS~= 0.000). This pattern
should be interpreted cautiously: the man-marking sample is small, and
zonal defences are likely deployed by teams facing higher-quality opponents,
introducing a selection confound. The primary SQ3 regression analysis
(Section~\ref{sec:sq3_analysis}) controls for this by operating on the full
1,401-corner dataset with defensive strategy labels from the automated
classifier.

\section{SQ4: Individual Player Qualities and Corner Success}
\label{sec:sq4}

\subsection{Corner Taker Quality}
\label{sec:sq4_taker}

The quality of the corner taker is characterised using three metrics.

\subsubsection{Delivery accuracy.}
For the 51~corners in the manually labelled set where the intended delivery zone
is known, delivery accuracy is defined as the proportion of corners landing in
the intended depth band (GA vs.\ PA vs.\ PE or deeper). This is reported as a
taker-specific statistic for takers with $\geq 5$~labelled deliveries.

\subsubsection{Delivery quality score (CTS-based).}
Taker quality is quantified as the \textit{Corner Threat Score} computed at the
taker level \citep{Lamberts2023}: the CTS formula (Equation~\ref{eq:cts}) is
applied to the taker's mean xA and mean sequence xG aggregated across all their
corners in the dataset:
\begin{equation}
  \mathrm{TakerCTS} = \frac{2 \cdot \overline{\mathrm{xG}}_{\mathrm{taker}}
                              \cdot \overline{\mathrm{xA}}_{\mathrm{taker}}}
                           {\overline{\mathrm{xG}}_{\mathrm{taker}}
                              + \overline{\mathrm{xA}}_{\mathrm{taker}}},
  \label{eq:taker_cts}
\end{equation}
where $\overline{\mathrm{xA}}_{\mathrm{taker}}$ is the taker's mean SciSports
expected assist per delivery and $\overline{\mathrm{xG}}_{\mathrm{taker}}$ is
the mean maximum shot xG in their corner sequences. TakerCTS integrates both
delivery quality and the threat generated by the taker's corners.
Only takers with $\geq 10$~corners in the dataset are assigned a score.

\subsubsection{Delivery consistency.}
The standard deviation of ball peak height ($\sigma_{z}$) across a taker's corners
is computed as a measure of trajectory consistency. Lower $\sigma_{z}$ indicates
greater repeatability of the delivery arc. Flight category entropy is reported
as a complementary measure.

\subsection{Target Player Quality}
\label{sec:sq4_target}

Where a shot or header immediately follows the corner delivery and the initiating
player can be identified from event data, that player is defined as the
\textit{aerial target}. Two quality indicators are extracted:

\begin{enumerate}
  \item \textbf{Aerial duel win rate}: the proportion of aerial duels won by the
        player across all match events (not limited to corners), computed from the
        full event dataset;
  \item \textbf{Physical attributes}: player height (metres), if available in the
        SciSports player metadata. Where not available, this metric is omitted and
        flagged as a data gap.
\end{enumerate}

\subsection{Controlling for Player Quality in SQ3}
\label{sec:sq4_control}

To isolate the effect of tactical setup (delivery type, defensive strategy) from
individual player contributions, the logistic regression in
Equation~\ref{eq:logit} is re-estimated after adding taker CTS and (where
available) target player height as additional covariates. A decrease in the
coefficients of tactical predictors after inclusion of player quality variables
would indicate that part of the observed tactical effectiveness is attributable
to individual quality rather than the tactic itself. This analysis addresses the
primary research question of SQ4.

\section{SQ5: Tactical Preparation Tool}
\label{sec:sq5}

\subsection{Overview}
\label{sec:sq5_overview}

SQ5 integrates the outputs of SQ1–SQ4 into a practical preparation tool aimed at
supporting coaches in designing corner kick routines against specific opponents.
The tool consists of two components: (A) an opponent scouting report and
(B) a positional response model, referred to informally as the
\textit{magnet board}.

\subsection{Component A: Opponent Scouting Report}
\label{sec:sq5_scouting}

For a given upcoming opponent, the following steps are executed:

\begin{enumerate}
  \item \textbf{Defensive profile}: all corners in which the opponent appeared
        as the defending team are identified and their defensive strategy types
        (from SQ1) are tabulated. The modal defensive strategy constitutes the
        \textit{typical defensive setup}.
  \item \textbf{Effectiveness lookup}: the SQ3 effectiveness matrix is queried
        for the opponent's typical defensive strategy. Offensive dimensions
        with above-average CTS against this strategy are highlighted as
        recommended attacking approaches.
  \item \textbf{Player qualities}: the opponent's man-markers (from SQ1) are
        identified, along with their aerial duel win rates (from SQ4), to
        estimate the risk of man-marking assignments.
\end{enumerate}

The scouting report is formatted as a standardised one-page summary
containing the opponent's defensive strategy distribution, the three most
effective offensive approaches (by mean xG), and a schematic of the opponent's
typical defensive organisation.

\subsection{Component B: Positional Response Model (Magnet Board)}
\label{sec:sq5_magnet}

The magnet board addresses the question: given a planned attacker configuration,
what defensive positioning is the opposing team likely to adopt, and where do
structural weaknesses appear?

\subsubsection{Representation.}
Each corner is represented by two spatial configurations, extracted from tracking
data at the delivery moment ($t = 0$):
\begin{itemize}
  \item \textbf{Attacker configuration}: the $(x, y)$ coordinates of all
        attacking outfield players (excluding the corner taker), recorded at
        $t - 3$~s (pre-delivery setup);
  \item \textbf{Defender configuration}: the $(x, y)$ coordinates of all
        defending outfield players at $t = 0$.
\end{itemize}

\subsubsection{Zonal defender positions.}
For teams with $\geq 10$~historical corners as the defending team, the mean
position of each defender classified as \texttt{ZONAL} is computed across all
their corners. This yields a \textit{zonal template}: a set of mean positions
for each zonal slot (typically 4–6 players), visualised as labelled markers
on a 2D pitch diagram. The zonal template is relatively stable across corners
for a given team, reflecting deliberate set-piece preparation; the standard
deviation of each position is reported to quantify consistency.

\subsubsection{Man-marker response.}
For defenders classified as \texttt{MAN}, the model estimates where a specific
man-marker is likely to position given an attacker placed at a given location.
Because man-markers track their assigned attacker, their expected position at
delivery is approximated by the attacker's delivery-moment position plus a
small offset vector derived from the mean positional lag observed in the
historical data. Formally, let $\mathbf{a}$ denote the planned attacker position
and $\hat{\boldsymbol{\delta}}$ denote the mean tracking offset:
\begin{equation}
  \hat{\mathbf{d}} = \mathbf{a} + \hat{\boldsymbol{\delta}},
  \label{eq:man_marker}
\end{equation}
where $\hat{\boldsymbol{\delta}}$ is estimated from the 95~man-marking pairs in
the manually labelled dataset. This offset captures the fact that man-markers
typically position slightly between their assigned attacker and the goal.

\subsubsection{Gap identification.}
Given a planned attacker configuration and the predicted defensive response
(zonal positions + man-marker responses), the uncovered pitch area is estimated
by computing, for each of the 12 central zones (GA1–PC3), the minimum distance
from any defensive player to the zone centroid. Zones with minimum distance
exceeding a threshold (calibrated on the labelled dataset) are flagged as
\textit{potentially open}, providing actionable guidance on which zone to target
with the delivery.

\subsubsection{Visualisation and interpretation.}
The magnet board is rendered as a standard 2D pitch diagram
(attacking half only) showing:
\begin{itemize}
  \item Zonal defenders as circles labelled with their role, at their expected
        mean position ($\pm$ one standard deviation ellipse);
  \item Man-markers as arrows from the planned attacker position to the
        predicted defender position;
  \item Potentially open zones highlighted in the target-zone grid.
\end{itemize}
The output is designed to be immediately legible to coaching staff without
statistical expertise. No machine learning model is used in the magnet board;
the tool is entirely rule-based and transparent.

\subsubsection{Limitation.}
The magnet board is trained on aggregate historical data and cannot account for
within-team variation or tactical adjustments in response to specific attacker
identities that the opposing coach has pre-planned. It should be used as a
probabilistic starting point for tactical discussion rather than a deterministic
prediction.

\section{Validation Strategy}
\label{sec:validation}

Model validation is summarised in Table~\ref{tab:validation}.

\begin{table}[h]
  \centering
  \caption{Validation approach for each analytical component.}
  \label{tab:validation}
  \begin{tabular}{p{4cm}p{4cm}p{4.5cm}}
    \toprule
    \textbf{Component} & \textbf{Ground truth} & \textbf{Validation method} \\
    \midrule
    Defender role model & 509 labelled defender rows & Leave-one-match-out CV \\
    Running line classifier & 858 labelled player rows & Leave-one-corner-out CV \\
    Attacker role classifier & 389 labelled player rows & Leave-one-corner-out CV \\
    Delivery swing & 51 labelled corners & Agreement rate \\
    Delivery zone & 51 labelled corners & Agreement rate \\
    Ball flight category & Manual check on 86 corners & Distribution validation \\
    Magnet board — zonal positions & 86 labelled corners & Mean position error (metres) \\
    Magnet board — man-marker offset & 95 labelled man-marking pairs & Mean offset error (metres) \\
    \bottomrule
  \end{tabular}
\end{table}

\section{Software and Reproducibility}
\label{sec:software}

All analyses were implemented in Python~3.10 using
\texttt{pandas}~\citep{pandas2020},
\texttt{numpy}~\citep{numpy2020},
\texttt{scikit-learn}~\citep{sklearn2011},
and \texttt{scipy}~\citep{scipy2020}.
Pitch visualisations were produced with \texttt{matplotlib}~\citep{matplotlib2007}.
Data extraction scripts, model training pipelines, and the magnet board
visualisation tool are version-controlled and available in the project repository.
All random seeds are fixed at 42 for reproducibility.

% ============================================================
%  BIBLIOGRAPHY ENTRIES (add to your .bib file)
%  \bibitem{Casal2015}
%    Casal, C.~A., Maneiro, R., Arill, T., Losada, J.~L., \& Ardá, A. (2015).
%    Analysis of corner kick success in elite football. \textit{International
%    Journal of Performance Analysis in Sport}, 15(2), 430--451.
%
%  \bibitem{Plakias2025}
%    Plakias, S., Armatas, V., \& Giakas, G. (2025). Technical–tactical analysis
%    of corner kicks in male soccer: A systematic review. \textit{Applied
%    Sciences}, 15(9), 4984.
%
%  \bibitem{Fernandez2019}
%    Fernández, J., \& Bornn, L. (2019). Wide open spaces: A statistical
%    technique for measuring space creation in professional soccer.
%    \textit{MIT Sloan Sports Analytics Conference}.
%
%  \bibitem{Lamberts2023}
%    Lamberts, M. (2023). Corner Threat Score: Measuring the threat a player
%    generates from corners. Medium / Analytics FC.
%
%  \bibitem{Rathke2017}
%    Rathke, A. (2017). An examination of expected goals and shot quality from a
%    soccer perspective. \textit{International Journal of Computer Science in
%    Sport}, 16(2), 76--96.
%
%  \bibitem{pandas2020} McKinney, W. (2010). Data Structures for Statistical
%    Computing in Python. \textit{Proceedings of the 9th Python in Science
%    Conference}.
%  \bibitem{numpy2020}  Harris, C.~R. et al. (2020). Array programming with
%    NumPy. \textit{Nature}, 585, 357--362.
%  \bibitem{sklearn2011} Pedregosa, F. et al. (2011). Scikit-learn: Machine
%    Learning in Python. \textit{JMLR}, 12, 2825--2830.
%  \bibitem{scipy2020}  Virtanen, P. et al. (2020). SciPy 1.0. \textit{Nature
%    Methods}, 17, 261--272.
%  \bibitem{matplotlib2007} Hunter, J.~D. (2007). Matplotlib: A 2D graphics
%    environment. \textit{Computing in Science \& Engineering}, 9(3), 90--95.
% ============================================================
